\chapter{Model Design and Implementation}
\section{Model Overview}
\subsection{Task Definition}
The objective of this work is to address the DTI prediction problem as a binary classification task.

\subsection{Model Inputs and Outputs}
The inputs and outputs were explained in detail in Chapter~\ref{chap:data-pipeline}, this section briefly goes over the inputs and outputs.
The inputs going into the model are Protein embeddings and Drug embeddings with strong semantic meaning.
The output is either a binder/non-binder classification determined by thresholds mentioned in Chapter~\ref{chap:data-pipeline}.

\subsection{Learning Paradigm}
The learning paradigm, how the model learns from the data, is supervised binary classification using PyTorch-based gradient optimisation using Adam \citep{kingma2015adam}.
The model learns to distinguish between binding and non-binding drug-target pairs based on labelled training data.

\section{Problem Formulation}
The prediction task is formulated as a supervised binary classification problem, where the goal is to conclude whether a given drug and protein pair produce a binding interaction.
The data instances of drug molecules and target proteins are represented as fixed length numerical vectors where $\mathbf{d} \in \mathbb{R}^{m}$ denotes the representation of a drug molecule and $\mathbf{p} \in \mathbb{R}^{n}$ denotes the representation of a target protein.

The goal of the model is to learn a function $f: \mathbb{R}^{m} \times \mathbb{R}^{n} \rightarrow [0,1],$ which maps a drug--protein pair $(\mathbf{d}, \mathbf{p})$ to a continuous interaction score $\hat{y}$ representing the predicted probability of binding:
$\hat{y} = f(\mathbf{d}, \mathbf{p}).$
To obtain a binary interaction label, the predicted score $\hat{y}$ is compared to the predefined decision threshold $\tau$, which is a tunable parameter; the value of $\tau$ = 0.3 is justified through data sensitivity analysis in Chapter~\ref{ch:-evaluation}:
\[
\hat{y}_{\text{class}} =
\begin{cases}
1, & \text{if } \hat{y} \ge \tau, \\
0, & \text{otherwise}.
\end{cases}
\]
The ground-truth labels $y \in \{0,1\}$ are derived from experimentally measured binding affinity values and binarised using affinity thresholds described in Chapter~\ref{chap:data-pipeline}.
Pairs exceeding the binding threshold are labelled as interacting, while pairs below the non-binding threshold are labelled as non-interacting.


Classification was selected over regression to improve robustness against experimental noise and to focus learning on more definitive binding interactions.
This formulation aligns with the intended application of the system, where the primary objective is to determine whether a candidate drug is likely to bind a given target rather than to predict an exact affinity value.

The learning objective is, therefore, to minimise the discrepancy between the predicted interaction labels $\hat{y}$ and the true labels $y$ in the training dataset, allowing the model to generalise to the unseen drug-target pairs.

\section{Integration of Input Representations}
Drug to Target interaction predictions require the combination of two different data styles, drugs being small chemical structures and proteins being amino acid sequences.
These representations differ in scale, semantic meaning, and dimensionality, making direct integration inappropriate.

To address this, the model uses a dual-input design, where drugs and proteins are processed independently through dedicated encoding pathways before being combined in a 128-dimensional space and fused for joint prediction.
Independent processing preserves the features and semantic meaning for each drug and protein, and joint learning is enabled once combined.
\subsection{Drug Representation}
Drug molecules are represented using MACCS fingerprints, which encode the presence or absence of predefined chemical substructures as a fixed length binary vector of 167 dimensions.

The fixed dimensionality makes it well suited for neural networks as it allows the model to learn patterns between specific substructures and binding behaviour.
This also makes it computationally efficient and reproducible across multiple experiments.
\subsection{Protein Representation}
Target proteins are represented using contextual embeddings created by the pretrained ProtBERT transformer model~\citep{Elnaggar2020.07.12.199554}.
Unlike fixed sequence-based embeddings, transformer-based embeddings capture bidirectional and context-dependent semantic meanings within amino acid sequences, capturing biochemical properties and structural information.

Each protein sequence is mapped to a fixed-length vector through mean pooling to produce a representation suitable for neural network processing; this is the aggregation strategy used in the original ProTrans work~\citep{Elnaggar2020.07.12.199554}.
Mean pooling provides a simple method for aggregating variable-length protein sequences while preserving global contextual information learned by the transformer.
This approach enables the model to leverage rich semantic information learned from large protein databases, improving generalisation to unseen targets.

\section{Model Architecture}
Due to the heterogeneous nature of dual inputs, the architecture is slightly modified from traditional neural networks.
Separate encoding pathways are used to process each input to a usable, unified representation for interaction prediction.

The \texttt{DrugEncoder} and \texttt{ProteinEncoder} are trainable multilayer perceptrons that project the precomputed MACCS and ProtBERT features into a shared 128-dimensional latent space.
Their parameters are updated jointly with the \texttt{InteractionPredictor} during training.
This two-stage design separates fixed, domain-informed feature extraction from learned task-specific adaptation: the representations contribute chemical and biological data, while the encoders learn to emphasise the most informative components of those representations for binding prediction.

\subsection{Drug Encoder}
The drug encoder is responsible for transforming molecular SMILES strings into a numerical representation using MACCS (Molecular ACCess System) keys \citep{durant2002reoptimization}.
The SMILES representation is parsed into an RDKit object, where MACCS fingerprint generation is applied, producing a 167-bit binary vector.

This representation captures essential structural properties of small molecules while maintaining a compact and consistent dimensionality.
In cases where a SMILES string cannot be successfully parsed, a zero vector is returned to preserve input dimensional consistency and prevent pipeline failure.

\subsection{Protein Representation}
Protein sequences are represented using contextual embeddings derived from a pretrained ProtBERT transformer model \citep{Elnaggar2020.07.12.199554}.
Each amino acid sequence is tokenised and passed through the frozen ProtBERT model to obtain per-token hidden states.
Mean pooling is then applied across the sequence dimension of the final hidden layer, producing a single 1024-dimensional embedding that summarises the biochemical and contextual properties of the entire sequence.

This process is computationally expensive, so each embedding is cached locally using hash-based lookup, reducing overhead when a sequence has already been computed.
This is particularly effective during dataset re-processing, as previously computed embeddings are loaded directly from the disk.

\subsection{Protein Encoder}
The protein encoder is a four-layer MLP (1024$\rightarrow$512$\rightarrow$256$\rightarrow$128) that compresses the precomputed ProtBERT embedding into a 128-dimensional latent representation.
This follows the design pattern used by~\citet{Singh2023}, where frozen pretrained protein language model embeddings are transformed into a lower-dimensional shared latent space by a learned nonlinear projection.
The layer widths being halved at each step, acts as a soft information bottleneck that encourages the encoder to only keep components of the ProtBERT embedding that are most informative for interaction prediction.
Unlike the frozen ProtBERT model, this encoder is trainable and participates in gradient-based optimisation during model training.

\subsection{Prediction Network}
\label{subsec:-prediction_network}
The prediction network contains three hidden layers (256, 128, 64) using ReLU activations for learning non-linear interactions between the drug and protein representations.
The network also makes use of dropout and L2 regularisation to prevent overfitting, which is a common issue in models with limited training data and high-dimensional inputs.
The sigmoid is applied only at inference, to map the raw logit to a probability.
During training, BCEWithLogitsLoss applies sigmoid internally when computing the loss for numerical stability.
The overall architecture is summarised in Table~\ref{tab:network_structure}.
\begin{table}[H]
\centering
\caption{Model Architecture and Training Hyperparameters}
\label{tab:network_structure}
\begin{tabular}{ll}
\toprule
\textbf{Parameter} & \textbf{Value} \\
\midrule
\textit{Architecture} \\
\quad Drug encoder         & 167$\rightarrow$256$\rightarrow$128 \\
\quad Protein encoder      & 1024$\rightarrow$512$\rightarrow$256$\rightarrow$128 \\
\quad Predictor            & 256$\rightarrow$256$\rightarrow$128$\rightarrow$64$\rightarrow$1 \\
\quad Activation           & ReLU \\
\quad Batch normalisation  & Yes \\
\quad Dropout              & 0.4 \\
\midrule
\textit{Training} \\
\quad Loss function        & BCEWithLogitsLoss \\
\quad Optimiser            & Adam \\
\quad Learning rate        & $1\times10^{-3}$ \\
\quad Weight decay         & $1\times10^{-4}$ \\
\quad Batch size           & 512 \\
\quad Max epochs           & 100 \\
\midrule
\textit{Regularisation} \\
\quad Early stopping       & Yes (patience = 15) \\
\quad LR scheduler         & CosineAnnealingLR ($T_{\max}=100$, $\eta_{\min}=1\times10^{-5}$) \\
\bottomrule
\end{tabular}
\end{table}

The full dual-branch architecture is defined as three PyTorch \texttt{nn.Module} classes: \texttt{DrugEncoder}, \texttt{ProteinEncoder}, and \texttt{InteractionPredictor}, composed inside the top-level \texttt{DTI\_DNN} module.
Rather than applying a uniform dropout rate throughout the network, a tapered schedule was used in which deeper layers receive lower dropout than earlier ones.
The reasoning is that deeper representations are more abstract and less redundant, so aggressive dropout at that stage risks destroying genuinely useful features rather than preventing the co-adaptation that dropout is designed to address.
A two-thirds multiplier ($\approx 0.67$) was chosen empirically: with a base rate of $0.4$, the deeper dropout layers therefore use approximately $0.27$ ($0.4 \times 0.67$).
Stronger tapers underused the regularisation budget in deeper layers, while weaker tapers produced less stable training; the $0.67$ value retained the benefit of regularisation without visible degradation during the iterative experiments described in Chapter~\ref{ch:-evaluation}.
The forward-pass definition of the top-level module is shown below.

\begin{lstlisting}[caption={dual-branch forward pass in DTI\_DNN}]
class DTI_DNN(nn.Module):
    def __init__(self, drug_input_dim=167, prot_input_dim=1024,
                 drug_hidden=256, prot_hiddens=(512, 256),
                 encoder_out=128, predictor_drop=0.4, encoder_drop=0.4):
        super().__init__()
        self.drug_encoder = DrugEncoder(drug_input_dim, drug_hidden,
                                        encoder_out, encoder_drop)
        self.prot_encoder = ProteinEncoder(prot_input_dim, prot_hiddens,
                                           encoder_out, encoder_drop)
        self.predictor = InteractionPredictor(encoder_out * 2, predictor_drop)

    def forward(self, x_drug, x_prot):
        drug_repr = self.drug_encoder(x_drug)
        prot_repr = self.prot_encoder(x_prot)
        fused = torch.cat([drug_repr, prot_repr], dim=1)
        return self.predictor(fused)
\end{lstlisting}

\subsection{Fusion and Prediction Layer}
The design choices reflected here were informed through iterative experimentation; the full progression is documented in Chapter~\ref{ch:-evaluation}.

Once both inputs have been encoded, both representations are combined to form a joint input for the prediction model.
The binary MACCS fingerprint with dimensionality $167$ and the ProtBERT embedding with dimensionality $1024$ are individually encoded to have dimensionality of $128$ before being concatenated into a single vector with dimensionality $256$ representing the drug-target pair.

This combined vector is then passed through the prediction network described in Table~\ref{tab:network_structure} and produces a binary classification based on the optimal threshold, $\tau = 0.3$, determined through experimentation.

\section{Training Procedure}
This section describes how the model parameters are optimised during training, an important process to how the model learns meaningful patterns from the training data.
\subsection{Loss Function}
The model is trained using Sigmoid + Binary Cross Entropy (BCE) loss, implemented via PyTorch's \texttt{BCEWithLogitsLoss}.
It is defined as:
\[
    E = -\frac{1}{N} \sum_{i=1}^{N} \left[ y_i \log(\hat{y}_i) + (1 - y_i) \log(1 - \hat{y}_i) \right],
\]
Where $N$ is the number of training samples, $y_i$ is the true binary label for the $i$-th sample, and $\hat{y}_i$ is the predicted probability of the positive class (binding interaction) for the $i$-th sample.
This loss function is more appropriate for binary classification because it penalises confident wrong predictions more heavily than MSE \citep{goodfellow2016deep}.

To address the class imbalance identified in Chapter~\ref{chap:data-pipeline} ($\sim$61\% binders, $\sim$39\% non-binders), the loss is weighted using PyTorch's \texttt{pos\_weight} argument.
This scales the positive-class contribution to the loss by the ratio of non-binders to binders computed on the training split, effectively giving underrepresented examples proportionally more influence during gradient updates without resampling the dataset.
\begin{lstlisting}[caption={class-imbalance weighting for BCEWithLogitsLoss}]
n_pos = y_train.sum().item()
n_neg = len(y_train) - n_pos
pos_weight = torch.tensor([n_neg / n_pos], dtype=torch.float32).to(device)
criterion = nn.BCEWithLogitsLoss(pos_weight=pos_weight)
\end{lstlisting}

\subsection{Optimisation Strategy}
\label{subsec:optimisation_strategy}
Model parameters were initially optimised using batch gradient descent with backpropagation.
Following iterative experimentation documented in Chapter~\ref{ch:-evaluation}, the model was reimplemented in PyTorch using the Adam optimiser with mini-batch gradient descent and He initialisation for weight initialisation.
Adam is an adaptive learning rate optimiser that maintains per-parameter learning rates using a combination of momentum and RMSProp, which can lead to faster convergence and improved performance compared to traditional stochastic gradient descent (SGD) \citep{kingma2015adam}.
Parameters are updated via:
\begin{equation}
    \theta_{t+1} = \theta_t - \frac{\alpha}{\sqrt{\hat{v}_t} + \epsilon} \hat{m}_t
\end{equation}
Where $\hat{m}_t$ and $\hat{v}_t$ are bias-corrected estimates for the first and second moments of the gradients, $\alpha$ is the learning rate, and $\epsilon$ is a small constant to prevent division by zero.
Following the default hyperparameters recommended by \citep{kingma2015adam}, this work uses $\alpha = 0.001$, $\beta_1 = 0.9$ $\beta_2 = 0.999$ and $\epsilon = 1 \times 10^{-8}$, with an L2 regularisation penalty applied via weight decay (see Table~\ref{tab:network_structure}).
Mini-batch gradient descent is used to balance the computational efficiency of batch gradient descent with the noise reduction benefits of stochastic gradient descent, allowing for more stable convergence and better generalisation \citep{masters2018revisiting}.
He initialisation is used to set the initial weights of the network, which helps to mitigate issues of vanishing or exploding gradients in deep networks, particularly when using ReLU activations \citep{he2015delving}.

\subsection{Hyperparameter Selection}
As shown in Table~\ref{tab:network_structure}, the configuration largely follows defaults set by Adam optimiser.
The learning rate was set to $1\times10^{-3}$, which is a common default for Adam and provides a good balance between convergence speed and stability.
Weight decay adds a penalty proportional to the squared magnitude of the weights to the loss function, which helps to prevent overfitting by adding a penalty on large weights.
This encourages the model to learn simpler patterns that generalise better to unseen data~\citep{Krogh1992}.
A value of $1\times10^{-4}$ was chosen through tuning, and it is consistent as a default value commonly used in PyTorch training pipelines.
The batch size of 512 was chosen to provide a good balance between computational efficiency and stable gradient estimation.
This sits at the upper boundary of the small-batch regime identified by~\citet{Keskar2017}, which is typically 32--512 and generalises better than larger batches.

\subsection{Regularisation}
The Dropout rate was determined through experimentation as seen in Chapter~\ref{ch:-evaluation}, and was set to 0.4.
This means that during training, 40\% of the neurons in the hidden layers are randomly deactivated in each iteration, helping to prevent overfitting by encouraging the model to learn more robust features that are not reliant on specific neurons~\citep{srivastava2014dropout}.
L2 regularisation is applied through the weight decay parameter in the Adam optimiser, which adds a penalty proportional to the square of the magnitude of the weights, encouraging smaller weights and thus simpler models that are less likely to overfit.

\section{Model Output Interpretation}
The model produces a single raw logit for each drug-target pair.
During training, \texttt{BCEWithLogitsLoss} applies the sigmoid function internally when computing the loss, avoiding numerical instability associated with a separate sigmoid layer.
At inference time, the sigmoid function is applied explicitly in the prediction service to map the logit to a probability $\hat{y} \in (0,1)$, representing the likelihood of binding (1) or non-binding (0).

To get a discrete prediction from the continuous output score at inference, the score is compared against a predefined threshold to optimise the performance of the model.
We settled on a threshold of 0.3 after experimentation as it produced the highest F1 score on the test set, which means it is the optimal balance between precision and recall.

The final prediction should be interpreted as a probability instead of a definitive prediction.
Therefore, this model is intended for the early-stage screening of drug-target pairs instead of replacing experimental validation.

\section{Implementation Details}
The model was originally implemented in Python using NumPy for the numerical computations.
All the forward and backward propagation steps were implemented manually for full transparency over the learning process and better insight into how changes affect the model's performance.
However, after conducting experiments and evaluating each improvement, the optimal implementation utilised PyTorch's framework.
The dual-branch design was a PyTorch-specific implementation using the \texttt{nn.Module}.



